%% conclusion.tex

\section{Conclusion}
\label{sec:conclusion}

\paragraph{Summary.}
This paper presented an empirical software engineering study of \rci{} in self-healing LLM agent frameworks, organized around three research questions that span prevalence, exploitability, and defense effectiveness. We analyzed the recovery paths of six production-grade frameworks (Reflexion, MetaGPT, Aider, OpenHands, AutoGen, LangGraph) at fixed commit hashes, constructed a recovery-pattern taxonomy using grounded theory coding, validated the identified structural pattern through controlled and end-to-end attack experiments on real LLM APIs, and evaluated the \trustorigin{} content-level provenance-tagging defense against naive baselines, stronger models, and adaptive adversaries.

\paragraph{Key findings.}
Our analysis yielded three principal findings, one per research question.

\textbf{RQ1 (Prevalence).} Five of six frameworks do not track the provenance of failure content on the recovery path: externally originated content (test output, file content, tool return values) is indistinguishable from internally generated diagnostics. Four of six frameworks (Reflexion, MetaGPT, Aider, AutoGen) exhibit the full \emph{trust-escalation trampoline} signature---the combination of no provenance tag and high-trust framing---which structurally upgrades untrusted content to trusted status on the recovery path. LangGraph provides partial provenance tagging (tool-exception status metadata), and OpenHands records error status in object metadata without reflecting it in the LLM-facing message framing.

\textbf{RQ2 (Exploitability and Severity).} The trust-escalation trampoline is practically exploitable and severe. The controlled testbed experiment on \system{}~\cite{omo} yields a baseline ASR of 0.57 (Wilson CI $[0.47, 0.66]$) with real LLM APIs. Cross-framework validation on LangGraph's public API confirms the vulnerability is not an artifact of the testbed (recovery-path ASR 0.76, real agentic-loop ASR 0.96). A confirmed end-to-end exploit on Reflexion's actual shipped retry path (commit \texttt{218cf0e}) achieves ASR 0.52 ($p = 3.5 \times 10^{-10}$ vs.\ 0\% benign false-positive rate), establishing exploitability on production code that real users run. The forward-defense bypass experiment shows \rci{} succeeds at ASR 0.68 even when three forward-path defenses are fully enabled ($p = 0.660$, indistinguishable from undefended), establishing \rci{} as an independent attack surface that forward-path defenses cannot cover.

\textbf{RQ3 (Defense Effectiveness).} The content-level \trustorigin{} tag reduces ASR from 0.57 to 0.07 ($p = 7.6 \times 10^{-15}$) on DeepSeek/Qwen, but degrades substantially on stronger instruction-following models: defended ASR rises to 0.50 on gpt-oss-120b and 0.36 on a closed-weight frontier model ($p = 0.027$). A defense-aware attacker raises defended ASR from 0.07 to 0.15, preserving an 84\% relative reduction versus baseline. The defense imposes 0\% false alarms and sub-millisecond machinery overhead. The central empirical finding is that \emph{content-level provenance defenses are effective but probabilistic, and degrade predictably as models become stronger instruction-followers}---quantifying the ceiling of content-level mitigation on the recovery channel.

\paragraph{Contributions.}
We contribute (i)~a recovery-pattern taxonomy classifying how self-healing frameworks handle failure-content provenance, with cross-case classification of six production frameworks; (ii)~a multi-layered exploitability evaluation establishing that the trampoline signature is practically exploitable on a controlled testbed, on a third-party framework (LangGraph), and on production code (Reflexion); (iii)~a forward-defense bypass experiment establishing \rci{} as an independent attack surface; (iv)~a defense-effectiveness evaluation quantifying the content-level defense's efficacy and its degradation profile under cross-model scaling and adaptive attack; and (v)~a statistical methodology for empirical security evaluation of LLM agent systems, combining Wilson 95\% confidence intervals, Fisher exact tests with Bonferroni correction where applicable, and Cohen's $\kappa$ inter-annotator agreement.

\paragraph{Implications.}
The findings carry three implications for the LLM agent framework community. First, the recovery path is a \emph{trust boundary}, not merely a robustness mechanism: provenance tracking at the failure-capture site is a necessary minimum, and high-trust framing applied uniformly to error content is a structural anti-pattern. Framework developers should audit recovery paths for the trampoline signature---the two binary criteria (Q1: provenance tag; Q2: high-trust frame) defined in \S\ref{sec:method-criteria} provide a lightweight screening test. Second, content-level defenses (untrusted framing, instruction hierarchy, spotlighting) provide probabilistic but not absolute protection: they degrade as models become stronger instruction-followers and are bypassable by adaptive adversaries. The empirical degradation profile quantified in RQ3 establishes a concrete ceiling for this defense class. Third, the recovery channel is an \emph{independent} attack surface that forward-path defenses cannot cover: input sanitizers, injection detectors, and tool-permission checks inspect the user's request, not the external content that the tool fetches and that causes the failure. Closing this gap requires defenses placed on the recovery path itself, not merely stronger forward-path defenses.

\paragraph{Future work.}
Three directions are most pressing. First, the single-round defense-aware attacker (\S\ref{sec:defense-aware}) provides a lower bound; a fully adaptive multi-round attacker with observed-rejection feedback could achieve higher ASR and should be evaluated. Second, the closed-weight frontier model evaluation (\S\ref{sec:closed-frontier}) covers one model accessed via its official API; extending to GPT-4o, Claude 3.5 Sonnet, and Gemini would strengthen the cross-model generalization claim. Third, the taxonomy (\S\ref{sec:rq1-taxonomy}) covers six frameworks at fixed commits; a longitudinal study tracking how provenance handling evolves over time would test whether the trampoline signature is converging toward or diverging from safe design as the ecosystem matures.
